\documentclass[11pt]{article}
\usepackage[a4paper, top=18mm, bottom=18mm, left=20mm, right=20mm]{geometry}
\usepackage[T2A]{fontenc}
\usepackage[utf8]{inputenc}
\usepackage[ukrainian]{babel}
\usepackage{graphicx}
\setlength{\parindent}{0pt}
\setlength{\parskip}{4pt}
\pagestyle{empty}
\begin{document}
%begin
{\large\textbf{Контрольна робота}}

\textbf{Варіант \No\,3}\hfill \textit{ПІБ:}\,\underline{\hspace{60mm}}
\vspace{4mm}

%beginitem
1. Що буде виведено за виконання фрагменту коду?
try:
    res = 8//8%8*8
    if isinstance(res, bool):
        print(f'bool;{res}')
    elif isinstance(res, int):
        print(f'int;{res}')
    elif isinstance(res, float):
        print(f'float;{res:.2f}')
    else:
        print('???')
except:
    print('Z')

%enditem

%beginitem
2. Відмітити все, що є коректним оператором С++:
( 1 ) n=1; n++

( 2 ) False=0

( 3 ) n=10; --n

( 4 ) x=float("inf")

%enditem

%beginitem
3. Відмітити все, що є однією правильною лексемою:
( 1 ) '123'

( 2 ) =<

( 3 ) fop

( 4 ) >

%enditem

%end
\newpage
\end{document}

